\section{Rounding with a small defect}\label{sec:rounding}

We need to round almost unimodular coefficients while controlling the
error on the whole circle.  The estimate below extends the real rounding
argument of \cite[Section~6]{OpenAILowerEnvelope2026} to prescribed complex
phases.  We quote its real matrix discrepancy lemma below; the proof
there uses the partial-coloring method of Spencer \cite{Spencer1985}
and the real-vector theorem of Lovett and Meka
\cite[Theorem~4 of arXiv:1203.5747v2]{LovettMeka2015}.

\begin{lemma}[Real matrix discrepancy]\label{lem:matrix-discrepancy}
There is an absolute constant $C$ such that, for integers $1\le s\le R$
and every $A\in[-1,1]^{R\times s}$, there is a vector
$\xi\in\{-1,1\}^s$ satisfying
\[
 \norm{A\xi}_\infty\le C\sqrt{s\log(2R/s)}.
\]
\end{lemma}

\begin{proof}
This is \cite[Lemma~6.2]{OpenAILowerEnvelope2026}, whose statement applies
to arbitrary real matrices in the displayed rectangular range.
\end{proof}

\begin{lemma}[Rounding with prescribed phases]\label{lem:complex-rounding}
Let $n\ge1$ be an integer and $Z_0,\ldots,Z_{n-1}\in\C$ satisfy
$|Z_j|\le1$.  Write $Z_j=d_j|Z_j|$, where $|d_j|=1$, choosing $d_j=1$
when $Z_j=0$, and put
\[
 \mu=\frac12\sum_{j=0}^{n-1}(1-|Z_j|).
\]
There are $\eta_j\in\{d_j,-d_j\}$ such that
\begin{equation}\label{eq:complex-rounding}
 \sup_{t\in\T}\left|\sum_{j=0}^{n-1}(\eta_j-Z_j)\e(jt)\right|
 \le C\left(1+\sqrt{\mu\log(80n/\mu)}\right),
\end{equation}
with an absolute constant $C$.  The square-root term is zero for
$\mu=0$, when the error can be zero.  For real inputs the outputs are
real signs.  The same estimate holds on any interval of $n$ consecutive
integer frequencies.
\end{lemma}

\begin{proof}
If $\mu=0$, take $\eta_j=Z_j$.  Suppose $\mu>0$, and set
$p_j=(1-|Z_j|)/2$.  Then $0\le p_j\le1/2$ and $\sum_jp_j=\mu$.
On the grid $t_\ell=\ell/(20n)$, $0\le\ell<20n$, form the real matrix
with $R=40n$ rows
\[
 A_{2\ell,j}=\Re\bigl(d_j\e(jt_\ell)\bigr),\qquad
 A_{2\ell+1,j}=\Im\bigl(d_j\e(jt_\ell)\bigr).
\]
Every entry belongs to $[-1,1]$.  We round $p$ to a vector in
$\{0,1\}^n$, preserving an upper bound on its total mass.

Choose an integer $J\ge1$ with $n2^{-J}\le1$, and first set
$p_j^{(J)}=2^{-J}\lfloor2^Jp_j\rfloor$.  This costs at most $1$ in
$\norm{A(p^{(J)}-p)}_\infty$ and does not increase the mass.
At stage $h=J,\ldots,1$, let $S_h$ contain those coordinates for which
$2^hp_j^{(h)}$ is odd, and let $s_h=|S_h|$.  If $s_h=0$, leave the
vector unchanged.  Otherwise apply Lemma~\ref{lem:matrix-discrepancy}
to the columns indexed by $S_h$.  Reverse all its signs if necessary so
that their sum is nonpositive, and put
\[
 p_j^{(h-1)}=
 \begin{cases}
 p_j^{(h)}+2^{-h}\xi_j,&j\in S_h,\\
 p_j^{(h)},&j\notin S_h.
 \end{cases}
\]
An odd multiple of $2^{-h}$ in $[0,1]$ lies between $2^{-h}$ and
$1-2^{-h}$.  Either change therefore stays in $[0,1]$ and lies on the
next coarser grid.  The total mass remains at most $\mu$, while the
simultaneous sign reversal leaves the absolute discrepancy unchanged.
In particular, each active coordinate has mass at least $2^{-h}$, so
\[
 s_h\le\min(n,2^h\mu),\qquad
 \norm{A(p^{(h-1)}-p^{(h)})}_\infty
 \le C2^{-h}\sqrt{s_h\log(80n/s_h)}.
\]
The last inequality is used only at nonempty stages.  Monotonicity of
$s\log(80n/s)$ on $(0,n]$ gives
\[
 s_h\log(80n/s_h)\le2^h\mu\log(80n/\mu).
\]
For $2^h\mu>n$, use the bound $n\log80$ on the left and
$\mu\le n/2$ on the right.  Summing the geometric series of step errors
yields, with $p'=p^{(0)}\in\{0,1\}^n$,
\begin{equation}\label{eq:rounding-matrix-error}
 \norm{A(p'-p)}_\infty
 \le1+C\sqrt{\mu\log(80n/\mu)}.
\end{equation}

Set $\eta_j=d_j(1-2p'_j)$ and
$Q(z)=\sum_{j=0}^{n-1}(\eta_j-Z_j)z^j$.
The real and imaginary parts of $Q(\e(t_\ell))$ are each bounded by
twice the left side of \eqref{eq:rounding-matrix-error}.  Thus the
claimed error bound holds on the grid.
To pass to the circle, let $S=\max_{|z|=1}|Q(z)|$.  The maximum principle,
applied also to the reversed polynomial, gives
\[
 |Q(z)|\le S\max(1,|z|^{n-1}),\qquad
 \max_{|z|\le1+1/n}|Q(z)|\le\exp(1)S.
\]
Cauchy's estimate in disks of radius $1/n$ centered on the unit circle
therefore gives
\[
 \left|\frac{d}{dt}Q(\e(t))\right|\le2\pi\exp(1)nS.
\]
Every point of $\T$ is within $1/(40n)$ of the grid.  Hence
$S\le\max_\ell|Q(\e(t_\ell))|+(\pi\exp(1)/20)S$.
Since $\pi\exp(1)/20<1$, this proves
\eqref{eq:complex-rounding}.  Multiplication by a unimodular monomial
handles shifted consecutive frequencies.  If the inputs are real,
all $d_j$ are signs and hence so are the outputs.
\end{proof}
